% ==================================================================
% Conscious Point Physics:
% The CPP Apparatus Model:
% Macroscopic Detectors, Decoherence, and the Quantum Processor
% Trade-Off Law
% CPP Companion SD#3 — Version 1
% GitHub: CPP/series_foundations/cpp_sd3_apparatus_model_v1.tex
% ==================================================================

\documentclass[11pt,a4paper]{article}
\usepackage[margin=1in]{geometry}
\usepackage[utf8]{inputenc}
\usepackage[T1]{fontenc}
\usepackage{amsmath,amssymb,amsthm}
\usepackage{graphicx}
\usepackage{caption,float,booktabs,parskip,xcolor,array}
\usepackage[numbers,sort&compress]{natbib}
\usepackage{hyperref}
\hypersetup{colorlinks=true,linkcolor=blue,urlcolor=cyan,citecolor=red}

\newtheorem{theorem}{Theorem}[section]
\newtheorem{proposition}[theorem]{Proposition}
\newtheorem{corollary}[theorem]{Corollary}
\newtheorem{remark}[theorem]{Remark}
\newtheorem{openprob}[theorem]{Open Problem}
\newtheorem{prediction}[theorem]{Prediction}

\newcommand{\lP}{l_P}
\newcommand{\tP}{t_P}
\newcommand{\SSV}{\mathrm{SSV}}
\newcommand{\hDP}{\mathrm{hDP}}
\newcommand{\eCP}{\mathrm{eCP}}
\newcommand{\phig}{\varphi}
\newcommand{\eps}{\varepsilon}

\title{\textbf{Conscious Point Physics:\\
The CPP Apparatus Model:\\
Macroscopic Detectors, Decoherence,\\
and the Quantum Processor Trade-Off Law}\\[6pt]
{\large Companion SD\#3 --- Version 1}}

\author{Thomas Lee Abshier, ND\\
Hyperphysics Institute\\
\url{https://hyperphysics.com}\\
\texttt{drthomas007@protonmail.com}}

\date{23 March 2026}

\begin{document}
\maketitle

\begin{abstract}
SD\#1~\citep{cpp_sd1} established that a measurement apparatus is a
CPP structure whose orientation $\theta$ couples to the hidden
variables $\lambda$ of the measured particle through the Nexus global
constraint.  SD\#2~\citep{cpp_sd2} derived the angular form of this
coupling from $H_4$ symmetry.  This paper derives the quantitative
aspects: the DP~Sea anisotropy imprinted by the apparatus, the
decoherence timescales, and the experimental trade-off governing
quantum processor tests of CPP.

Four results are established:

\textbf{Theorem~1 (Classical apparatus):}
A macroscopic detector with $N_{\rm app} \sim 10^{26}$ atoms at
temperature $T$ decoheres in time $\tau_{\rm dec} = \hbar /
(N_{\rm app} k_B T) \lesssim 10^{-39}$~s.  The apparatus is
always classical on any experimentally accessible timescale.
This is why quantum mechanics has worked so well for 100 years
--- not because deviations are impossible, but because macroscopic
detectors are never quantum.

\textbf{Theorem~2 (Quantum processor trade-off):}
For a quantum processor with $N_q$ qubits at temperature $T$,
the Nexus participation fraction $\eps_{\rm QP} = N_q / N_{\rm sub}$
and the quantum coherence time $\tau_q = \hbar/(N_q k_B T)$
satisfy:
\begin{equation*}
\eps_{\rm QP} \cdot \tau_q
= \frac{\hbar}{N_{\rm sub} k_B T} = \text{const}(T).
\end{equation*}
Increasing $N_q$ raises $\eps_{\rm QP}$ but proportionally shortens
$\tau_q$.  The \emph{signal per circuit run} is independent of $N_q$.

\textbf{Theorem~3 (DP~Sea anisotropy):}
An oriented apparatus of linear size $L$ imprints a fractional
anisotropy $\delta = \text{sea\_strength} \cdot (\lP/L) \sim 3
\times 10^{-34}$ on the surrounding DP~Sea.  This anisotropy is
the CPP microscopic description of what ``apparatus orientation
$\theta$'' means.

\textbf{Prediction~1 (N-qubit scaling law):}
For $N_q$ parallel Bell pairs on a quantum processor, the total
accumulated CPP signal scales linearly with $N_q$ (because $N_q$
parallel pairs each contribute an independent CPP correction).
The functional signature is $\delta E \propto N_q$, while noise
grows as $\sqrt{N_q}$.  The signal-to-noise ratio for CPP detection
scales as $\sqrt{N_q}$, making large-qubit processors the
optimal testing ground.
\end{abstract}

\tableofcontents
\newpage

%=====================================================================
\section{Introduction}
%=====================================================================

SD\#1 and SD\#2 established the \emph{form} of the CPP deviation
from standard Bell correlations.  This paper addresses the
\emph{scale}: how large is the deviation, why is it undetectable in
current experiments, and what experimental architecture gives the
best chance of detection?

The answers require a CPP model of the measurement apparatus.  In
CPP, an apparatus is not an external classical device --- it is a
collection of $\sim 10^{26}$ eCPs in the DP~Sea, participating in
the same Nexus constraint as the measured particles.  Its orientation
$\theta$ is a collective degree of freedom of those eCPs.  Its
thermal decoherence rate determines whether it behaves classically
or quantum-mechanically.

The central result is a conservation law (Theorem~2): the product
$\eps_{\rm QP} \cdot \tau_q$ is a constant at fixed temperature.
This has the structure of an uncertainty principle for CPP
detectability: you can have high Nexus sensitivity (large $\eps$)
or long quantum coherence (large $\tau$), but not both.  The
solution is not more qubits but more \emph{repetitions}: running
many short circuits accumulates CPP signal additively while noise
only accumulates as $\sqrt{M}$.

%=====================================================================
\section{The Apparatus as a CPP Structure}
\label{sec:apparatus_cpp}
%=====================================================================

\subsection{Collective orientation in the DP~Sea}

A polarisation analyser (or spin detector) is, in CPP, a
macroscopic collection of $N_{\rm app}$ eCPs arranged in a
crystal, magnetic, or optical medium.  Its orientation $\theta$
is the macroscopic average of the individual eCP SSV vectors:
\begin{equation}
\theta = \frac{1}{N_{\rm app}}\sum_{i=1}^{N_{\rm app}}\theta_i,
\label{eq:theta_collective}
\end{equation}
where $\theta_i$ is the orientation of the $i$-th eCP relative to
a reference direction.  Individual eCPs fluctuate thermally; the
collective average is sharp to $O(1/\sqrt{N_{\rm app}})$.

\subsection{DP~Sea anisotropy from the apparatus}

An oriented collection of eCPs creates an anisotropic SSV field
in the surrounding DP~Sea.  The fractional anisotropy at position
$\mathbf{r}$ relative to the apparatus centre is:
\begin{equation}
\frac{\delta\rho_{\rm DP}(\mathbf{r},\theta)}{\rho_0}
\equiv \delta(\mathbf{r},\theta)
= \text{sea\_strength} \cdot \frac{\lP}{|\mathbf{r}|}
  \cdot \cos(\theta - \phi_{\mathbf{r}}),
\label{eq:dp_anisotropy}
\end{equation}
where $\phi_{\mathbf{r}}$ is the direction angle of $\mathbf{r}$
from the apparatus centre and $\rho_0$ is the isotropic DP~Sea
density.

\begin{theorem}[DP~Sea anisotropy scale]
\label{thm:anisotropy}
At the surface of the apparatus (characteristic size $L$), the
fractional DP~Sea anisotropy is:
\begin{equation}
\delta_0 = \text{sea\_strength} \cdot \frac{\lP}{L}
\approx 0.185 \times \frac{1.62\times10^{-35}\,\text{m}}{L}.
\label{eq:delta0}
\end{equation}
For a 1~cm apparatus: $\delta_0 \approx 3 \times 10^{-34}$.
This is the CPP microscopic meaning of ``the apparatus has
orientation $\theta$'': the DP~Sea surrounding the apparatus
is anisotropic by the fraction $\delta_0$ in the direction $\theta$.
\end{theorem}

\begin{proof}
The SSV field of a single eCP is proportional to sea\_strength
times $\hbar c / r^2$ (in natural units where $\lP$ sets the
scale).  The collective SSV field of $N_{\rm app}$ coherently
aligned eCPs scales as $N_{\rm app} \times \lP/r^2$.
At the apparatus surface $r = L$, normalising by $\rho_0$ and
taking the leading dipole term gives~\eqref{eq:delta0}.\qed
\end{proof}

\begin{remark}[Why $\delta_0$ is relevant to the Nexus]
The Nexus operates atemporally across the entire lattice.  The
DP~Sea anisotropy $\delta_0$ is the \emph{causal} imprint of the
apparatus on the surrounding DP~Sea --- the part that propagates
via the SSV at $c$.  The Nexus contribution (SD\#1
Proposition~1~\citep{cpp_sd1}) is additional and atemporal,
scaling as $\eps_{\rm Nexus} = N_{\rm part}/N_{\rm app}$.  The
two contributions are independent channels by which the apparatus
orientation influences the particle's DP~Sea state.
\end{remark}

%=====================================================================
\section{Thermal Decoherence of the Apparatus}
\label{sec:decoherence_app}
%=====================================================================

\begin{theorem}[Apparatus is always classical]
\label{thm:classical}
A macroscopic measurement apparatus with $N_{\rm app}$ atoms at
temperature $T$ decoheres to a classical state in time:
\begin{equation}
\tau_{\rm dec}^{\rm app} = \frac{\hbar}{N_{\rm app}\,k_B T}.
\label{eq:tau_app}
\end{equation}
For all experimentally accessible conditions,
$\tau_{\rm dec}^{\rm app} \ll \tau_{\rm meas}$, so the apparatus
is always in a classical (definite-orientation) state during any
Bell measurement.
\end{theorem}

\begin{proof}
The decoherence time~\eqref{eq:tau_app} follows from the standard
Joos--Zeh--Zurek theory of quantum decoherence applied to a
macroscopic object \citep{joos1985,zurek1991}: a superposition of
two macroscopically distinct orientations decoheres when the
thermal fluctuation energy $N_{\rm app} k_B T$ exceeds the quantum
of uncertainty $\hbar$ divided by the coherence time.  The CPP
derivation follows by applying the Lindblad master equation of
companion 2040d~\citep{cpp2040d} to the collective orientation
degree of freedom.\qed
\end{proof}

\begin{table}[H]
\centering
\caption{Decoherence time~\eqref{eq:tau_app} for a macroscopic
  apparatus ($N_{\rm app} = 10^{26}$ atoms) at various temperatures.
  All values are far shorter than any measurement timescale
  ($\tau_{\rm meas} \gtrsim 10^{-9}$~s).}
\label{tab:tau_app}
\begin{tabular}{@{}lcc@{}}
\toprule
Condition & $T$ & $\tau_{\rm dec}^{\rm app}$ \\
\midrule
Room temperature & $300$~K & $2.5\times10^{-40}$~s \\
Liquid nitrogen & $77$~K & $9.9\times10^{-40}$~s \\
Liquid helium & $4.2$~K & $1.8\times10^{-38}$~s \\
Dilution refrigerator & $10$~mK & $7.6\times10^{-36}$~s \\
\bottomrule
\end{tabular}
\end{table}

\begin{remark}[Why classical detectors do not suppress CPP]
The classical character of the apparatus does \emph{not} mean the
CPP correction vanishes.  It means the apparatus orientation $\theta$
is sharp and well-defined during the measurement.  The Nexus
correlation~(SD\#1, eq.~3) couples $\theta$ to the particle hidden
variable $\lambda$; this coupling is active regardless of whether
the apparatus is in a quantum superposition.  The CPP correction
exists \emph{because} the apparatus has a definite orientation,
not despite it.
\end{remark}

%=====================================================================
\section{The Quantum Processor: Coherence vs.\ Participation}
\label{sec:QP}
%=====================================================================

\subsection{The two roles of qubit count}

A quantum processor with $N_q$ qubits at temperature $T$ (dilution
refrigerator, $T \sim 10$~mK) plays two roles in CPP Bell tests:

\begin{enumerate}
\item \textbf{Coherence time:} The qubits form a quantum system with
  collective decoherence time $\tau_q = \hbar/(N_q k_B T)$, which
  shrinks as $N_q$ increases.
\item \textbf{Participation fraction:} With $N_q$ qubits, the
  particle-side of the Nexus constraint involves $N_q$ Grid~Points,
  raising the participation fraction to $\eps_{\rm QP} = N_q /
  N_{\rm sub}$, which grows as $N_q$ increases.
\end{enumerate}

These two effects are in direct competition.

\begin{theorem}[Quantum processor trade-off law]
\label{thm:tradeoff}
For a quantum processor with $N_q$ qubits at temperature $T$,
the product of the Nexus participation fraction $\eps_{\rm QP}$
and the quantum coherence time $\tau_q$ is:
\begin{equation}
\eps_{\rm QP} \cdot \tau_q
= \frac{N_q}{N_{\rm sub}} \cdot \frac{\hbar}{N_q k_B T}
= \frac{\hbar}{N_{\rm sub}\, k_B T}
= \mathrm{const}(T).
\label{eq:tradeoff}
\end{equation}
The product is independent of $N_q$ and depends only on temperature
and the substrate size $N_{\rm sub}$.
\end{theorem}

\begin{proof}
Direct substitution of $\eps_{\rm QP} = N_q/N_{\rm sub}$ and
$\tau_q = \hbar/(N_q k_B T)$.\qed
\end{proof}

\begin{remark}[Uncertainty-principle analogy]
Equation~\eqref{eq:tradeoff} has the structure of an uncertainty
principle: increasing $N_q$ raises sensitivity ($\eps_{\rm QP}
\uparrow$) but shortens the coherence window ($\tau_q \downarrow$)
by exactly the same factor.  There is no optimal $N_q$ in terms of
the product; instead, the relevant figure of merit is the total
signal accumulated over many circuit runs (see
Prediction~\ref{pred:scaling}).
\end{remark}

\subsection{Signal per circuit run}

The CPP signal accumulated in a single circuit run of $N_{\rm
gates}$ gates is:
\begin{equation}
\mathcal{S}_{\rm run}
= \eps_{\rm QP} \cdot |f_{H_4}(\theta)| \cdot N_{\rm gates}
= \frac{N_q}{N_{\rm sub}} \cdot |f_{H_4}|
  \cdot \frac{\tau_q}{\tau_{\rm gate}},
\label{eq:signal_run}
\end{equation}
where $\tau_{\rm gate} \approx 50$~ns is the gate time and
$N_{\rm gates} = \tau_q / \tau_{\rm gate}$.  Substituting
Theorem~\ref{thm:tradeoff}:
\begin{equation}
\mathcal{S}_{\rm run}
= \frac{\hbar}{N_{\rm sub}\, k_B T \cdot \tau_{\rm gate}}
  \cdot |f_{H_4}(\theta)|.
\label{eq:signal_independent}
\end{equation}

\begin{corollary}[Signal per run is independent of $N_q$]
\label{cor:Nq_independent}
The CPP signal accumulated in a single maximum-depth circuit run is
independent of the number of qubits $N_q$.  Numerically, at
$T = 10$~mK and $\tau_{\rm gate} = 50$~ns:
\begin{equation}
\mathcal{S}_{\rm run} \approx
\frac{1.055\times10^{-34}}{10^{26} \times 1.38\times10^{-23}
\times 10^{-2} \times 5\times10^{-8}}
\approx 1.5 \times 10^{-28}.
\label{eq:Srun_numeric}
\end{equation}
\end{corollary}

\begin{remark}[What $N_q$ actually buys]
Although $N_q$ does not affect the signal per run, it does help
in a different way.  With $N_q$ qubits, one can measure $N_q$
parallel Bell pairs \emph{simultaneously} in a single run.  Each
parallel pair contributes an independent CPP correction.  The
total signal from $N_q$ parallel pairs per run scales as $N_q$,
while the noise (shot noise) scales as $\sqrt{N_q}$.  The
signal-to-noise ratio per run therefore scales as $\sqrt{N_q}$
--- larger processors are better, but by a square root, not linearly.
\end{remark}

\subsection{The $N_q$-scaling law}

\begin{prediction}[Quantum processor signal-to-noise scaling]
\label{pred:scaling}
For a Bell test on $N_q$ simultaneous qubit pairs, run $M$ times:
\begin{align}
\text{Total CPP signal} &\sim N_q \cdot M \cdot \mathcal{S}_{\rm run}
\propto N_q M, \label{eq:total_signal} \\
\text{Total noise (shot)} &\sim \sqrt{N_q \cdot M}, \label{eq:total_noise} \\
\text{SNR}_{\rm CPP} &\sim \sqrt{N_q \cdot M} \cdot \mathcal{S}_{\rm run}.
\label{eq:SNR}
\end{align}
A CPP signal, if present, will be identified by its \emph{linear
growth with $N_q$} at fixed $M$, distinguishing it from:
\begin{itemize}
\item Systematic errors (independent of $N_q$).
\item Shot noise (scales as $\sqrt{N_q}$, not $N_q$).
\item Other quantum effects (would require theoretical
  identification of why they scale with $N_q$).
\end{itemize}
\end{prediction}

%=====================================================================
\section{Temperature Dependence}
\label{sec:temperature}
%=====================================================================

From Theorem~\ref{thm:tradeoff}, $\eps_{\rm QP} \cdot \tau_q
\propto 1/T$.  Colder systems accumulate CPP signal faster.
The signal per run~\eqref{eq:signal_independent} scales as $1/T$:
\begin{equation}
\mathcal{S}_{\rm run}(T) = \frac{\hbar}{N_{\rm sub}\, k_B T
  \cdot \tau_{\rm gate}} \propto \frac{1}{T}.
\label{eq:S_T}
\end{equation}

\begin{table}[H]
\centering
\caption{Signal per circuit run~\eqref{eq:signal_independent} vs.\
  temperature, for $N_{\rm sub} = 10^{26}$ and
  $\tau_{\rm gate} = 50$~ns.  Cooling from 10~mK to 1~mK gives
  a factor of 10 enhancement.}
\label{tab:signal_T}
\begin{tabular}{@{}lcc@{}}
\toprule
Temperature & $\eps_{\rm QP}\cdot\tau_q$ & $\mathcal{S}_{\rm run}$
  (per run, per pair) \\
\midrule
$100$~mK & $7.6\times10^{-37}$~s & $1.5\times10^{-29}$ \\
$10$~mK (current) & $7.6\times10^{-36}$~s & $1.5\times10^{-28}$ \\
$1$~mK & $7.6\times10^{-35}$~s & $1.5\times10^{-27}$ \\
$0.1$~mK & $7.6\times10^{-34}$~s & $1.5\times10^{-26}$ \\
\bottomrule
\end{tabular}
\end{table}

\begin{prediction}[Temperature enhancement]
\label{pred:temperature}
At fixed $N_q$ and $M$, the CPP signal scales as $1/T$.
Dilution refrigerators at $0.1$~mK (achievable but uncommon)
give a factor of 100 over current $10$~mK systems.
Distinguishing CPP from temperature-independent systematic errors
is possible by running the same circuit at multiple temperatures:
a genuine CPP signal should show the $1/T$ dependence of~\eqref{eq:S_T}.
\end{prediction}

%=====================================================================
\section{The Nexus Coupling $\Gamma(\theta, \lambda)$}
\label{sec:nexus_coupling}
%=====================================================================

Combining the DP~Sea anisotropy model (Theorem~\ref{thm:anisotropy})
and the Nexus participation fraction (SD\#1, Proposition~1), the
total coupling of the apparatus orientation $\theta$ to the particle
hidden variable $\lambda$ is:
\begin{equation}
\Gamma(\theta, \lambda)
= \underbrace{\eps_{\rm Nexus} \cdot f_{H_4}(\theta - \theta_\lambda)}
             _{\text{Nexus (atemporal)}}
+ \underbrace{\delta_0 \cdot \cos(\theta - \theta_\lambda)}
             _{\text{direct SSV (causal)}},
\label{eq:Gamma_total}
\end{equation}
where $\theta_\lambda$ is the preferred axis of the particle's
hidden DP~Sea configuration and $f_{H_4}$ is the angular function
derived in SD\#2.

The direct SSV contribution (second term) falls as $\lP/L \sim
10^{-34}$ and is negligible.  The Nexus contribution (first term)
at $\eps_{\rm Nexus} \sim 10^{-26}$ dominates by eight orders of
magnitude, confirming that the CPP effect on Bell correlations is
dominated by the atemporal Nexus constraint, not by the causal
SSV field.

\begin{remark}[The Nexus term vs.\ the DP~Sea term]
This hierarchy is physically significant.  If the CPP correction
arose from the direct SSV field (causal), it would propagate at the
speed of light and the angular dependence would be $\cos\theta$ ---
indistinguishable from a systematic error in the standard Bell
correlation.  Because the Nexus term dominates, the angular
dependence is $f_{H_4}(\theta)$ with its 5-fold and 3-fold
structure (SD\#2).  The icosahedral angular signature is the
distinctive fingerprint that separates CPP from any causal correction.
\end{remark}

%=====================================================================
\section{Why Standard Bell Tests Are Quantitatively Blind}
\label{sec:why_blind}
%=====================================================================

Combining the magnitude estimate (SD\#1) with the apparatus model:

\begin{table}[H]
\centering
\caption{Why current Bell tests cannot detect CPP, and what
  would be needed.  $\mathcal{S}_{\rm run}$ is the signal per
  circuit run from Corollary~\ref{cor:Nq_independent}.}
\label{tab:sensitivity}
\begin{tabular}{@{}p{4.5cm}p{3cm}p{4.5cm}@{}}
\toprule
Experiment & Sensitivity & CPP signal \\
\midrule
Standard optical Bell test & $\sim 10^{-3}$ & $\sim 10^{-26}$ \\
Loophole-free Bell (2015) & $\sim 10^{-3}$ & $\sim 10^{-26}$ \\
Best precision Bell (2018) & $\sim 10^{-3}$ & $\sim 10^{-26}$ \\
Near-future Bell ($10^8$ runs) & $\sim 10^{-6}$ & $\sim 10^{-26}$ \\
\midrule
10-qubit QP Bell, $10^6$ runs & $\sim 10^{-5}$ & $\sim 10^{-27}$ \\
1000-qubit QP Bell, $10^9$ runs & $\sim 10^{-7}$ & $\sim 10^{-25}$ \\
Hypothetical $T=0.1$~mK system & $\sim 10^{-8}$ & $\sim 10^{-26}$ \\
\bottomrule
\end{tabular}
\end{table}

The gap between sensitivity and signal is currently $\sim 23$
orders of magnitude.  No foreseeable technology bridges this gap
directly.  However, Hossenfelder's prediction~\citep{hossenfelder2024}
is that quantum computer engineering will produce \emph{unexpected
anomalies} --- not from directly targeting CPP, but from building
ever-larger and more precise quantum systems.  The CPP prediction is
that those anomalies, when they come, will have the specific angular
and scaling signature of Predictions~\ref{pred:scaling}
and~\ref{pred:temperature}.

%=====================================================================
\section{Consolidated Open Problems}
%=====================================================================

\begin{openprob}[Derive $\delta_0$ from the SSV field equations]
\label{op:delta0}
Theorem~\ref{thm:anisotropy} gives the scaling $\delta_0 \sim
\text{sea\_strength} \cdot \lP/L$ from the SSV dipole field.
A rigorous derivation from the CPP SSV field equations (SR/GR
companion C2~\citep{abshier_sr}) would include the numerical
prefactor and the full angular dependence of
$\delta(\mathbf{r},\theta)$ for a realistic polariser geometry
(linear vs.\ circular vs.\ crystal-face polariser).
\end{openprob}

\begin{openprob}[Decoherence rate from the CPP Lindblad equation]
\label{op:decoherence_rate}
Theorem~\ref{thm:classical} quotes the standard
Joos--Zeh--Zurek decoherence time.  A CPP-specific derivation
using the Lindblad master equation of companion 2040d~\citep{cpp2040d}
would include corrections from the DP~Sea structure (the DP~Sea
has a granular Planck-scale structure that standard quantum
field theory ignores).  These corrections are $O((\lP/\lambda_{\rm
dB})^2) \sim 10^{-58}$ for room-temperature objects and completely
negligible, but the derivation would close the logical gap.
\end{openprob}

\begin{openprob}[Nexus coupling beyond the participation fraction]
\label{op:nexus_coupling}
The Nexus participation fraction $\eps_{\rm Nexus} = N_{\rm part}
/ N_{\rm app}$ is the leading-order estimate of the Nexus coupling
strength (SD\#1 Proposition~1~\citep{cpp_sd1}).  A rigorous
derivation from the Nexus constraint $\sum_i \delta b_i = 0$
requires integrating the constraint over the apparatus degrees of
freedom explicitly.  This is the core calculation of SD\#4.
\end{openprob}

\begin{openprob}[Multi-qubit entanglement and the Nexus]
\label{op:multiqubit}
Prediction~\ref{pred:scaling} assumes that $N_q$ qubit pairs each
contribute independent CPP corrections.  If the qubits are
\emph{mutually entangled} (not just pairwise), the hidden variable
$\lambda$ spans all $N_q$ qubits simultaneously.  The Nexus
coupling between an $N_q$-qubit entangled state and the $N_{\rm
sub}$-atom substrate may scale differently --- potentially as
$N_q^2/N_{\rm sub}$ rather than $N_q/N_{\rm sub}$.  Deriving the
correct scaling for multi-qubit entangled states is an open problem.
\end{openprob}

%=====================================================================
\section{Experimental Programme Summary}
%=====================================================================

Combining SD\#1--3, the complete experimental programme for testing
CPP is:

\begin{table}[H]
\centering
\caption{The three-paper experimental prescription from SD\#1--3.}
\label{tab:programme}
\begin{tabular}{@{}p{2cm}p{5cm}p{5.5cm}@{}}
\toprule
From & What it says & How to test \\
\midrule
SD\#1 & Magnitude $\eps \sim N_{\rm part}/N_{\rm app} \sim 10^{-26}$ &
  Baseline; current experiments 22 orders short \\[4pt]
SD\#2 & Angular form: $f_{H_4}(\theta)$ with extrema at $\{36°, 72°, 120°\}$;
  null at $90°$ &
  Scan $E(\theta)$ over $[0°, 180°]$; standard CHSH (45°) is not optimal;
  use ratio test at H$_4$-special angles \\[4pt]
SD\#3 & Signal $\propto N_q$ (parallel pairs), $\propto 1/T$;
  trade-off law $\eps \cdot \tau = \text{const}$ &
  Use largest available $N_q$; cool to lowest $T$;
  run many repetitions; check linear $N_q$ and $1/T$ scalings \\
\bottomrule
\end{tabular}
\end{table}

The full CPP signature in a quantum processor Bell test is:
\begin{enumerate}
\item A residual $\delta E(\theta)$ after subtracting $-\cos\theta$
  with 72° and 120° periodicity (SD\#2).
\item The residual grows linearly with $N_q$ (SD\#3
  Prediction~\ref{pred:scaling}).
\item The residual grows as $1/T$ (SD\#3 Prediction~\ref{pred:temperature}).
\item The ratio $\delta E(36°)/\delta E(120°) \approx -1.065$ is
  $\phig$-valued and independent of amplitude (SD\#2 eq.~21).
\item The residual vanishes at $\theta = 90°$ to leading order (SD\#2).
\end{enumerate}

All five properties together would uniquely identify a CPP signal.
Properties 2--4 can distinguish CPP from any systematic error with
a standard angular dependence.

%=====================================================================
\section{Conclusion}
%=====================================================================

The CPP apparatus model establishes three quantitative results.

The apparatus decoherence theorem (Theorem~\ref{thm:classical})
confirms that macroscopic detectors are always classical ---
consistent with 100 years of Bell tests showing no anomalies.  The
absence of detected CPP deviations is not a problem for the theory;
it is what is expected given the $\sim 10^{-26}$ magnitude of the
correction.

The quantum processor trade-off law (Theorem~\ref{thm:tradeoff})
shows that the product $\eps_{\rm QP} \cdot \tau_q = \hbar /
(N_{\rm sub} k_B T)$ is a constant at fixed temperature.  More
qubits give higher sensitivity but proportionally shorter coherence.
The net CPP signal per circuit run is independent of $N_q$
(Corollary~\ref{cor:Nq_independent}) but the SNR for $N_q$ parallel
Bell pairs scales as $\sqrt{N_q}$.

The DP~Sea anisotropy (Theorem~\ref{thm:anisotropy}) gives the
microscopic CPP meaning of apparatus orientation: a fractional
$\sim 3 \times 10^{-34}$ anisotropy in the surrounding DP~Sea.
The Nexus term dominates over the causal SSV term by eight orders
of magnitude, so the angular signature is icosahedral ($f_{H_4}$
from SD\#2), not cosine-shaped.

The three scaling predictions --- linear $N_q$, $1/T$ temperature
dependence, and the 72°/120° angular periodicity --- together provide
a complete, falsifiable experimental signature for CPP quantum
foundations.

\bibliographystyle{unsrtnat}
\begin{thebibliography}{14}

\bibitem{cpp_sd1}
T.~L. Abshier,
``The Nexus as superdeterministic mechanism,''
CPP SD\#1 v1, Hyperphysics Institute (2026).

\bibitem{cpp_sd2}
T.~L. Abshier,
``The 600-cell angular structure of Bell deviations,''
CPP SD\#2 v1, Hyperphysics Institute (2026).

\bibitem{cpp_qm_series}
T.~L. Abshier and Grok,
``Quantum mechanics from the 600-cell lattice,''
CPP QM Series 2040a--f v3.1, Hyperphysics Institute (2026).

\bibitem{cpp2040d}
T.~L. Abshier and Grok,
``Open quantum systems and the Lindblad equation,''
CPP-2040d v3.1, Hyperphysics Institute (2026).

\bibitem{cpp_ew_unified}
T.~L. Abshier and Grok,
``The electroweak sector from the 600-cell lattice,''
CPP EW Unified v2.1, Hyperphysics Institute (2026).

\bibitem{hossenfelder2024}
S.~Hossenfelder,
``Why faster than light travel is possible after all,''
YouTube lecture transcript (2024).

\bibitem{hossenfelder_palmer}
S.~Hossenfelder and T.~Palmer,
``Rethinking superdeterminism,''
\textit{Front.\ Phys.} \textbf{8}, 139 (2020).

\bibitem{joos1985}
E.~Joos and H.~D. Zeh,
``The emergence of classical properties through interaction with
the environment,''
\textit{Z.\ Phys.\ B} \textbf{59}, 223 (1985).

\bibitem{zurek1991}
W.~H. Zurek,
``Decoherence and the transition from quantum to classical,''
\textit{Phys.\ Today} \textbf{44}(10), 36 (1991).

\bibitem{bell1964}
J.~S. Bell,
``On the Einstein--Podolsky--Rosen paradox,''
\textit{Physics} \textbf{1}, 195 (1964).

\bibitem{li2018}
M.-H. Li \textit{et al.},
``Test of local realism into the past,''
\textit{Phys.\ Rev.\ Lett.} \textbf{121}, 080404 (2018).

\bibitem{abshier_sr}
T.~L. Abshier,
``Mechanistic derivation of relativistic effects via Space Stress
Vector,'' viXra:2511.0062 (2025).

\bibitem{brouwer1989}
A.~E. Brouwer, A.~M. Cohen, and A.~Neumaier,
\textit{Distance-Regular Graphs}, Springer (1989).

\bibitem{pdg2026}
Particle Data Group,
``Review of Particle Physics 2026,''
\textit{Phys.\ Rev.\ D} \textbf{110}, 030001 (2026).

\end{thebibliography}

\end{document}
